% TODO: Complete the Implementation chapter

% \subsection{Development Environment}

% % TODO: Describe your development environment
% % Include: OS, IDEs, version control, etc.

% \textbf{[TODO: Add development environment details here]}

% \subsection{Implementation Overview}

% % TODO: Provide an overview of the implementation process
% % What was the development methodology?
% % What were the main phases?

% \textbf{[TODO: Add implementation overview here]}

% \subsection{Core Components}

% % TODO: Describe the implementation of core components
% % For each major component:
% % - Explain its purpose
% % - Describe key implementation details
% % - Mention any challenges faced and how they were resolved


% ============================================================================
% 4.1 DEVELOPMENT ENVIRONMENT
% ============================================================================

\subsection{Development Environment}

\subsubsection{Hardware Configuration}

The entire project was developed and tested on the following hardware:

\begin{itemize}
    \item \textbf{Processor:} Apple M4 with 10-core CPU (4 performance cores + 6 efficiency cores)
    \item \textbf{Memory:} 16 GB unified memory
    \item \textbf{GPU:} Apple M4 GPU with 10 cores supporting Metal 3
    \item \textbf{Storage:} 512 GB SSD
    \item \textbf{Operating System:} macOS Sonoma (14.0+)
\end{itemize}

\subsubsection{Software Configuration}

The development environment consisted of:

\begin{itemize}
    \item \textbf{Programming Language:} Python 3.14
    \item \textbf{IDE:} Visual Studio Code with Python extension, Jupyter Notebook 7.0
    \item \textbf{Version Control:} Git (2.40+) and GitHub for source code management
    \item \textbf{Virtual Environment:} venv for dependency isolation
    \item \textbf{Package Management:} pip (24.0+) for installing dependencies
    \item \textbf{Testing Framework:} pytest (8.0+) for unit testing
\end{itemize}

\subsubsection{Core Libraries and Dependencies}

\begin{table}[H]
\centering
\caption{Core libraries and their purposes}
\label{tab:libraries}
\begin{tabularx}{\textwidth}{|l|l|X|}
\hline
\textbf{Library} & \textbf{Version} & \textbf{Purpose} \\
\hline
NumPy & 2.4.2 & Efficient numerical computations and matrix operations \\
\hline
Matplotlib & 3.10.0 & Generating comprehensive visualizations and plots \\
\hline
Pillow (PIL) & 10.0.0 & Image loading, processing, and manipulation \\
\hline
Scikit-learn & 1.5.0 & K-Means clustering implementation and metrics \\
\hline
JAX & 0.9.1 & GPU-accelerated computations on Apple M4 \\
\hline
PyTorch & 2.5.0 & Metal Performance Shaders acceleration \\
\hline
SciPy & 1.17.1 & Scientific computing and signal processing \\
\hline
pytest & 8.0.0 & Unit testing framework \\
\hline
\end{tabularx}
\end{table}

% ============================================================================
% 4.2 IMPLEMENTATION OVERVIEW
% ============================================================================

\subsection{Implementation Overview}

\subsubsection{Development Methodology}

The project followed an **iterative, test-driven development (TDD)** methodology. The implementation process consisted of four main phases:

\begin{enumerate}
    \item \textbf{Phase 1: Foundation (Weeks 1-2)} – Set up project structure, implemented base classes, and established the common interface for all compression algorithms.
    
    \item \textbf{Phase 2: Core Algorithms (Weeks 3-6)} – Implemented entropy coding methods (Huffman, Arithmetic, LZW) and spatial domain techniques (RLE, variable block size).
    
    \item \textbf{Phase 3: Advanced Algorithms (Weeks 7-10)} – Developed predictive coding (JPEG lossless), context-based methods (CALIC, LOCO-I), and transform domain techniques.
    
    \item \textbf{Phase 4: AI/ML Methods (Weeks 11-14)} – Implemented clustering (K-Means, Gath-Geva), optimization (PSO), and neural network (FBPNN) approaches with GPU acceleration.
\end{enumerate}

\subsubsection{Project Structure}

The codebase was organized into a modular, layered architecture:

\begin{verbatim}
project/
├── data/               # Image I/O and preprocessing
├── algorithms/         # All compression algorithm implementations
│   ├── entropy/        # Huffman, Arithmetic, LZW
│   ├── spatial/        # RLE, variable block, lossy+residual
│   ├── predictive/     # JPEG lossless prediction
│   ├── context/        # CALIC, LOCO-I
│   ├── transform/      # S-Transform, S+P, lifting wavelet
│   ├── clustering/     # K-Means, Gath-Geva
│   ├── optimization/   # PSO
│   └── neural/         # FBPNN
├── metrics/            # Compression ratio, MSE, PSNR
├── visualization/      # Plotting and figure generation
├── utils/              # Helper functions (bit I/O, entropy)
├── tests/              # Unit tests for all algorithms
└── notebooks/          # Jupyter notebooks for demonstrations
\end{verbatim}

\subsubsection{Common Interface Design}

All compression algorithms implement a consistent interface:

\begin{verbatim}
class BaseCompressor:
    """Base class for all compression algorithms"""
    
    def compress(self, image):
        """
        Compress an image.
        
        Args:
            image: numpy array of shape (H, W) or (H, W, C)
            
        Returns:
            compressed_data: Compressed representation
            metadata: Dictionary containing compression metadata
        """
        pass
    
    def decompress(self, compressed_data, metadata):
        """
        Decompress to reconstruct original image.
        
        Returns:
            image: Reconstructed numpy array
        """
        pass
    
    def evaluate(self, original, reconstructed):
        """Calculate MSE, PSNR, and compression ratio"""
        pass
\end{verbatim}

% ============================================================================
% 4.3 CORE COMPONENTS
% ============================================================================

\subsection{Core Components}

\subsubsection{Entropy Coding Components}

\textbf{Purpose:} To exploit symbol frequency imbalances for lossless compression.

\textbf{Key Implementation Details:}
\begin{itemize}
    \item \textbf{Huffman Coding}: Built a binary tree using a min-heap priority queue. Generated prefix-free codes where frequent symbols get shorter codes.
    \item \textbf{Arithmetic Coding}: Used integer arithmetic with 32-bit precision to avoid floating-point errors. Implemented cumulative frequency tables for efficient encoding.
    \item \textbf{LZW}: Maintained a dictionary of previously seen substrings; dynamically grew the dictionary during compression.
\end{itemize}

\textbf{Challenges and Solutions:}
\begin{itemize}
    \item \textit{Challenge:} Arithmetic coding precision issues.
    \item \textit{Solution:} Implemented integer-based arithmetic coding instead of floating-point.
\end{itemize}

\subsubsection{GPU-Accelerated Components}

\textbf{Purpose:} To accelerate computationally intensive algorithms (Gath-Geva, PSO, FBPNN) using Apple M4 Metal.

\textbf{Key Implementation Details:}
\begin{itemize}
    \item Used JAX with Metal backend for GPU acceleration
    \item Applied JIT compilation using \texttt{@jit} decorator
    \item Implemented vectorized operations with \texttt{vmap} for batch processing
    \item Added automatic fallback to CPU when GPU is unavailable
\end{itemize}

\begin{verbatim}
@jit
def gath_geva_update(X, U, m, c):
    """GPU-accelerated Gath-Geva update step"""
    Um = U ** m
    numerator = jnp.dot(Um.T, X)
    denominator = jnp.sum(Um, axis=0).reshape(-1, 1)
    return numerator / (denominator + 1e-10)
\end{verbatim}

\textbf{Challenges and Solutions:}
\begin{itemize}
    \item \textit{Challenge:} Singular covariance matrices in Gath-Geva.
    \item \textit{Solution:} Added regularization term (\(1\times10^{-6}I\)) to covariance matrices.
\end{itemize}

\subsubsection{Visualization Components}

\textbf{Purpose:} To generate comprehensive visualizations for algorithm comparison.

\textbf{Key Implementation Details:}
\begin{itemize}
    \item Side-by-side comparison of original and compressed images
    \item Error maps using hot colormap (brighter = more error)
    \item Histograms of pixel values and residuals
    \item Convergence plots for iterative algorithms
    \item Color palettes for clustering methods
\end{itemize}

\textbf{Challenges and Solutions:}
\begin{itemize}
    \item \textit{Challenge:} Visualizing large error ranges.
    \item \textit{Solution:} Used logarithmic scaling and adaptive colormap limits.
\end{itemize}

\subsubsection{Testing Components}

\textbf{Purpose:} To ensure correctness and lossless reconstruction.

\textbf{Key Implementation Details:}
\begin{itemize}
    \item Unit tests for each algorithm using pytest
    \item Test images included uniform, gradient, random, and natural images
    \item Verified lossless reconstruction for all lossless algorithms
    \item Benchmarked execution time for performance evaluation
\end{itemize}

\begin{verbatim}
def test_huffman_lossless():
    original = np.random.randint(0, 256, (64, 64), dtype=np.uint8)
    compressed, metadata = huffman.compress(original)
    reconstructed = huffman.decompress(compressed, metadata)
    assert np.array_equal(original, reconstructed)
\end{verbatim}

\textbf{Challenges and Solutions:}
\begin{itemize}
    \item \textit{Challenge:} Ensuring consistent results across different platforms.
    \item \textit{Solution:} Fixed random seeds and used deterministic algorithms.
\end{itemize}

\subsubsection{Optimization Techniques}

The following optimization techniques were employed throughout the implementation:

\begin{itemize}
    \item \textbf{Vectorization}: Used NumPy vectorized operations instead of Python loops (10-50x speedup)
    \item \textbf{JIT Compilation}: Applied JAX's just-in-time compilation for critical functions (5-10x speedup)
    \item \textbf{Memory Efficiency}: Implemented lazy loading and streaming for large images
    \item \textbf{Batch Processing}: Processed pixels in batches for clustering algorithms
    \item \textbf{Caching}: Cached intermediate results to avoid recomputation
\end{itemize}

\subsection{Installation and Usage}

The project can be installed and run using the following commands:

\begin{verbatim}
# Clone the repository
git clone https://github.com/faizanrefai/IITJ-MTP-Template-Generator.git
cd IITJ-MTP-Template-Generator

# Create and activate virtual environment
python -m venv venv
source venv/bin/activate  # On macOS/Linux

# Install dependencies
pip install -r requirements.txt

# Run tests
pytest tests/

# Run the demo notebook
jupyter notebook notebooks/demo.ipynb
\end{verbatim}

\subsubsection{Huffman Coding}


\subsubsection{Source Entropy and the Shannon Bound}

Let an image be represented as a discrete source that emits pixel intensity values from an alphabet $\mathcal{A} = \{0, 1, \dots, M-1\}$, where $M = 256$ for 8-bit grayscale images. For a given pixel value $i \in \mathcal{A}$, let $p_i$ denote its empirical probability (i.e., frequency divided by the total number of pixels). The source entropy $H(I)$ is defined as:

\begin{equation}
    H(I) = -\sum_{i=0}^{M-1} p_i \log_{2} p_i \quad \text{(bits per pixel)}
\end{equation}

This quantity represents the theoretical lower bound on the average number of bits per pixel required to encode the image losslessly \cite{VemuriSectionII}. The practical implication is that no lossless compression algorithm can achieve an average bit rate lower than $H(I)$; the degree to which an algorithm approaches this bound serves as a primary measure of its coding efficiency.

\paragraph{Compression Ratio and Bit Rate}

The performance of a lossless compression algorithm is typically quantified by the \textit{compression ratio} ($CR$) or the \textit{bit rate} ($BR$, measured in bits per pixel). For an image of dimensions $R \times C$ originally stored with $b$ bits per pixel, the original size is $R \cdot C \cdot b$ bits. If the compressed representation occupies $S$ bits, then:

\begin{equation}
    CR = \frac{R \cdot C \cdot b}{S}, \quad BR = \frac{S}{R \cdot C}
\end{equation}

The two metrics are related by the simple identity:

\begin{equation}
    CR = \frac{b}{BR}
\end{equation}

A compression ratio greater than $1$ indicates effective compression, while a ratio less than $1$ signifies expansion. The bit rate $BR$ should ideally be close to the source entropy $H(I)$ for an efficient algorithm.


\subsubsection{Component 2}
\textbf{[TODO: Add implementation details for component 2]}

\subsubsection{Component 3}
\textbf{[TODO: Add implementation details for component 3]}

\subsection{Key Features}

% TODO: Describe the implementation of key features
% Include screenshots or code snippets where appropriate

\textbf{[TODO: Add key features implementation here]}

\subsection{Database Implementation}

% TODO: Describe database schema and implementation
% Include ER diagram or schema diagram if helpful

\textbf{[TODO: Add database implementation details here]}

\subsection{API Implementation}

% TODO: Describe API endpoints and their implementation
% Include API documentation or examples

\textbf{[TODO: Add API implementation details here]}

\subsection{User Interface}

% TODO: Describe UI implementation
% Include screenshots of key screens

\textbf{[TODO: Add UI implementation details here]}

\subsection{Testing}

% TODO: Describe testing approach and implementation
% Include:
% - Unit testing
% - Integration testing
% - User testing
% - Test coverage

\textbf{[TODO: Add testing details here]}

\subsection{Challenges and Solutions}

% TODO: Discuss major challenges faced during implementation
% How were they resolved?

Table~\ref{tab:challenges} summarizes the main challenges and their solutions.

\begin{table}[H]
\centering
\caption{Implementation Challenges and Solutions}
\label{tab:challenges}
\begin{tabular}{@{}p{5cm}p{8cm}@{}}
\toprule
\textbf{Challenge} & \textbf{Solution} \\ \midrule
TODO: Challenge 1 & TODO: Solution 1 \\
TODO: Challenge 2 & TODO: Solution 2 \\
TODO: Challenge 3 & TODO: Solution 3 \\
\end{tabular}
\end{table}